% ==============================================================================
% Chapter 4: Implementation and Testing
% ==============================================================================
\chapter{Implementation and Testing}
\label{ch:impl_testing}

This chapter documents the step-by-step software implementation of the proposed multi-agent traffic signal system, the physical and virtual experimental environments, the rigorous functional verification test cases, and the safety measures integrated into the software build.

\section{Development Details}
\label{sec:ch4_dev_details}

The development workflow was executed sequentially across seven structured development steps, matching the work breakdown structure defined in Chapter 1.

\subsection*{Step 1: Environment Setup and Dependency Configuration}
The project was built in an isolated Python 3.10 virtual environment (\texttt{.venv}) configured with PyTorch 2.2 for GPU/CPU tensor operations, Ultralytics YOLOv8 for real-time computer vision inference, OpenCV 5.0 for video stream rasterization, and Streamlit for web telemetry. Listing \ref{lst:setup_cmds} details the shell commands used to establish the development environment and verify GPU/CPU execution availability.

\begin{lstlisting}[language=bash, caption={Environment initialization and dependency installation commands}, label={lst:setup_cmds}]
# Create and activate Python 3.10 virtual environment
python -m venv .venv
.venv\Scripts\activate

# Install core deep learning, vision, and simulation dependencies
pip install torch torchvision --index-url https://download.pytorch.org/whl/cu121
pip install ultralytics opencv-python gymnasium streamlit plotly pandas numpy

# Verify hardware accelerator acceleration and dependency versions
python -c "import torch, cv2, ultralytics; print(f'PyTorch: {torch.__version__} | CUDA: {torch.cuda.is_available()} | OpenCV: {cv2.__version__}')"
\end{lstlisting}

\subsection*{Step 2: Empirical Dataset Ingestion and Transition Tuple Synthesis}
To anchor the system to authentic Indian road behavior, an automated ingestion pipeline was developed in \texttt{tools/download\_and\_prep\_datasets.py}. The script ingests:
\begin{enumerate}
    \item \textbf{IIT Delhi Traffic Density Dataset}: 4 days (September 12--15) comprising 216,004 records of continuous 1-second interval queue and stop density data from 6 approach cameras across 3 signalized intersections in Delhi NCR.
    \item \textbf{Mendeley Pune Heterogeneous Count Dataset (\texttt{xnf2k6n288/1})}: Empirical turning counts across 3 urban junctions, establishing the baseline modal split: 58\% Two-Wheelers, 21\% Passenger Cars, 14\% Auto-Rickshaws, 4\% City Buses, and 3\% Heavy Trucks, yielding an empirical conversion factor of $1 \text{ vehicle} \equiv 0.850 \text{ PCU}$.
    \item \textbf{IISc AIM BMD-45 Taxonomy}: 14-class Indian vehicle mapping with IRC:106 PCU weights stored in \texttt{data/processed/bmd45\_taxonomy.json}.
\end{enumerate}
The pipeline synthesized \textbf{72,001 full transition tuples} $(s_t, a_t, r_t, s_{t+1})$, serialized into a 5.36 MB compressed archive at \texttt{data/processed/delhi\_pune\_transitions.npz}.

\subsection*{Step 3: Neural Dynamics World Model Architecture and Training}
In \texttt{src/neural\_world\_model.py}, the \texttt{TrafficDynamicsWorldModel} was implemented using residual linear blocks, Layer Normalization, and multi-task prediction heads for state differential $\Delta s_t \in \mathbb{R}^{66}$ and immediate scalar reward $r_t \in \mathbb{R}$. Listing \ref{lst:world_model_code} presents the core forward pass implementation.

\begin{lstlisting}[language=Python, caption={Residual Multi-Task Neural World Model architecture (\texttt{src/neural\_world\_model.py})}, label={lst:world_model_code}]
class TrafficDynamicsWorldModel(nn.Module):
    def __init__(self, state_dim=66, action_dim=9, hidden_dim=256):
        super().__init__()
        self.input_proj = nn.Linear(state_dim + action_dim, hidden_dim)
        
        # Residual Blocks with LayerNorm and GELU
        self.block1 = nn.Sequential(
            nn.LayerNorm(hidden_dim),
            nn.Linear(hidden_dim, hidden_dim),
            nn.GELU(),
            nn.Linear(hidden_dim, hidden_dim)
        )
        self.block2 = nn.Sequential(
            nn.LayerNorm(hidden_dim),
            nn.Linear(hidden_dim, hidden_dim),
            nn.GELU(),
            nn.Linear(hidden_dim, hidden_dim)
        )
        
        # Multi-task heads: State differential and reward
        self.state_head = nn.Linear(hidden_dim, state_dim)
        self.reward_head = nn.Linear(hidden_dim, 1)

    def forward(self, state, action_onehot):
        x = torch.cat([state, action_onehot], dim=-1)
        h = self.input_proj(x)
        h = h + self.block1(h)
        h = h + self.block2(h)
        
        delta_state = self.state_head(h)
        reward = self.reward_head(h)
        return delta_state, reward.squeeze(-1)
\end{lstlisting}

The world model was trained via \texttt{tools/train\_world\_model.py} for 15 epochs on 61,200 training transitions with an AdamW optimizer ($\text{lr} = 10^{-3}$, weight decay $= 10^{-4}$) and Cosine Annealing learning rate schedule.

\subsection*{Step 4: Data-Driven Gymnasium Environment Implementation}
To achieve complete simulator independence, \texttt{src/data\_driven\_env.py} was constructed as a 100\% drop-in replacement for the standard Gymnasium \texttt{MAPPOEnv}. It embeds the trained \texttt{TrafficDynamicsWorldModel}, tracks approach queues across 8 incoming lanes per junction, scales arrival demand by operational scenarios, and enforces the dual-ring signal safety state machine.

\subsection*{Step 5: Multi-Agent PPO Policy Training}
The MAPPO training script was implemented in \texttt{src/train\_mappo\_data\_driven.py}. The training pipeline initialized 4 parallel vectorized instances of \texttt{DataDrivenTrafficEnv}, executing 30,000 corridor timesteps across 30 policy updates. Listing \ref{lst:actor_critic} details the Actor policy and Centralized Critic implementations in \texttt{src/mappo\_networks.py}.

\begin{lstlisting}[language=Python, caption={Actor and Centralized Critic networks (\texttt{src/mappo\_networks.py})}, label={lst:actor_critic}]
class Actor(nn.Module):
    def __init__(self, obs_dim=22, n_actions=3, hidden_dim=128):
        super().__init__()
        self.net = nn.Sequential(
            nn.Linear(obs_dim, hidden_dim),
            nn.Tanh(),
            nn.Linear(hidden_dim, hidden_dim),
            nn.Tanh(),
            nn.Linear(hidden_dim, n_actions)
        )
        # Orthogonal initialization
        for m in self.net.modules():
            if isinstance(m, nn.Linear):
                nn.init.orthogonal_(m.weight, gain=np.sqrt(2))
                nn.init.constant_(m.bias, 0.0)

    def get_action(self, obs, deterministic=False):
        logits = self.net(obs)
        dist = torch.distributions.Categorical(logits=logits)
        action = torch.argmax(logits, dim=-1) if deterministic else dist.sample()
        return action, dist.log_prob(action), dist.entropy()

class Critic(nn.Module):
    def __init__(self, state_dim=66, hidden_dim=256):
        super().__init__()
        self.net = nn.Sequential(
            nn.Linear(state_dim, hidden_dim),
            nn.Tanh(),
            nn.Linear(hidden_dim, hidden_dim),
            nn.Tanh(),
            nn.Linear(hidden_dim, 1)
        )
\end{lstlisting}

\subsection*{Step 6: Edge Vision Perception Engine}
In \texttt{vehicle\_detection/live\_vision\_pipeline.py}, the perception engine wraps YOLOv8 Nano (\texttt{yolov8n.pt}) with multi-object tracking (Linear Assignment Problem) and Indian vehicle classification. It extracts approach bounding boxes, calculates instantaneous PCU loads, and builds the normalized 22-dimensional observation vector $o_i$ consumed by the trained Actor policy.

\subsection*{Step 7: Dual-Interface Live Deployment Tools}
Two complementary interfaces were implemented:
\begin{enumerate}
    \item \textbf{High-Speed OpenCV Video HUD (\texttt{src/live\_traffic\_controller.py})}: A standalone 30+ FPS video application that overlays bounding boxes, active PCU metrics, and an animated 3-intersection signal status bar directly on camera frames.
    \item \textbf{Streamlit Interactive Web Dashboard (\texttt{src/dashboard\_live\_vision.py})}: A browser-based monitoring suite rendering live camera streams, interactive Plotly charts, 3-color signal pods (Red, Yellow, Green), and live demonstration triggers.
\end{enumerate}

\subsection{Implementation Challenges and Engineering Resolutions}
Table \ref{tab:ch4_challenges} documents the key software challenges encountered during system construction and the technical measures adopted to resolve them.

\begin{table}[htbp]
\centering
\small
\caption{Implementation challenges and technical resolutions}
\label{tab:ch4_challenges}
\begin{tabularx}{\textwidth}{|p{3.5cm}|c|X|}
\hline
\textbf{Technical Challenge} & \textbf{Impact} & \textbf{Resolution Implemented} \\ \hline
TraCI Socket Disconnects & Critical & Replaced the SUMO socket interface entirely with a pre-trained PyTorch Deep Residual Neural World Model (\texttt{data\_driven\_env.py}). \\ \hline
Windows Console Unicode Encoding Errors & Medium & Console prints threw \texttt{UnicodeEncodeError} when printing Unicode block characters; replaced with standard ASCII bars (\texttt{\#}, \texttt{-}). \\ \hline
Streamlit Duplicate Element IDs in Streaming Loop & High & Calling \texttt{plotly\_chart} repeatedly in an infinite loop caused \texttt{StreamlitDuplicateElementId}; replaced with native \texttt{st.bar\_chart()} and \texttt{st.line\_chart()} with throttled placeholder updates. \\ \hline
Signals Never Showing Red Light & High & Modified \texttt{DataDrivenTrafficEnv} and dashboard state machines to strictly enforce 3-second Yellow clearance followed by 3-second All-Red clearance. \\ \hline
\end{tabularx}
\end{table}

\section{Experimental Setup}
\label{sec:ch4_exp_setup}
The physical hardware, operating systems, and software dependencies used to benchmark the prototype are documented in Table \ref{tab:ch4_hardware_setup}.

\begin{table}[htbp]
\centering
\small
\caption{Hardware and software experimental configuration}
\label{tab:ch4_hardware_setup}
\begin{tabularx}{\textwidth}{|p{4.2cm}|X|}
\hline
\textbf{Component} & \textbf{Specification / Configuration} \\ \hline
\textbf{Host Processor} & Intel Core i5-11400H @ 2.70 GHz (6 Cores, 12 Threads) / AMD Ryzen 5 \\ \hline
\textbf{System Memory} & 16 GB DDR4 3200 MHz RAM \\ \hline
\textbf{Hardware Accelerator} & NVIDIA GeForce RTX 3050 Laptop GPU (4 GB VRAM) / Intel UHD Graphics \\ \hline
\textbf{Operating System} & Microsoft Windows 11 Home (64-bit, Build 22631) \\ \hline
\textbf{Python Runtime} & Python 3.10.11 in isolated virtual environment (\texttt{.venv}) \\ \hline
\textbf{Deep Learning Stack} & PyTorch 2.2.2+cu121, Torchvision 0.17.2, NumPy 2.2.6 \\ \hline
\textbf{Vision Engine} & Ultralytics 8.4.156, OpenCV-Python 5.0.0, Lap 0.5.14 \\ \hline
\textbf{Web Telemetry} & Streamlit 1.64.0, Plotly 5.20.0, Pandas 2.2.1 \\ \hline
\textbf{Test Video Datasets} & \texttt{intersection\_video.mp4} ($1152 \times 720$ @ 30 FPS, 1582 frames), \texttt{synthetic\_indian\_traffic.mp4} ($1280 \times 720$ @ 30 FPS) \\ \hline
\end{tabularx}
\end{table}

\section{Test Cases and Functional Verification}
\label{sec:ch4_test_cases}

\subsection{Verification Strategy}
Verification was conducted across three distinct testing tiers:
\begin{enumerate}
    \item \textbf{Unit Component Verification}: Individual validation of the \texttt{TrafficDynamicsWorldModel} forward pass, the Actor logit generation, the PCU mapping logic, and the safety state machine transitions.
    \item \textbf{Integration Testing}: End-to-end verification passing synthetic and recorded video frames through YOLO detection, observation vector synthesis, Actor inference, and HUD rendering.
    \item \textbf{System Resilience Testing}: Air-gapped execution with physical network adapters disabled, verifying zero cloud or internet dependencies.
\end{enumerate}

\subsection{Functional Test Cases}
Table \ref{tab:ch4_test_cases} details the six formal test cases used to validate the operational integrity of the system.

\begin{table}[htbp]
\centering
\scriptsize
\setlength{\tabcolsep}{3pt}
\caption{System test cases, expected behavior, and observed outcomes}
\label{tab:ch4_test_cases}
\begin{tabularx}{\textwidth}{|c|p{2.2cm}|p{2.6cm}|p{2.8cm}|p{3.2cm}|c|}
\hline
\textbf{ID} & \textbf{Scenario} & \textbf{Input / Condition} & \textbf{Expected Behavior} & \textbf{Observed Output} & \textbf{Status} \\ \hline
\textbf{TC-01} & Indian Vehicle Classification & Video frame with bikes, autos, cars, buses & Correctly classifies 5 classes and assigns IRC:106 PCU weights & Detected 2W (0.5 PCU), Auto (1.0 PCU), Car (1.0 PCU), Truck (3.0 PCU) & \textbf{PASS} \\ \hline
\textbf{TC-02} & Stationary Queue Filtering & Vehicles waiting at red stop line & Flags vehicles moving $< 4$ px over 4 frames as stopped & Halting vehicles flagged; queue PCU correctly aggregated & \textbf{PASS} \\ \hline
\textbf{TC-03} & World Model Prediction & Global state (66-dim) + Joint action (9-dim) & Predicts $\Delta \hat{s}_t$ with MSE $< 0.005$; runtime $< 2$ ms & State MSE $= 0.001579$; Inference latency $= 0.42$ ms & \textbf{PASS} \\ \hline
\textbf{TC-04} & Multi-Agent Actor Decision & Local 22-dim observation vector $o_i$ & Outputs discrete action $a_i \in \{0, 1, 2\}$ in $< 5$ ms & Optimal phase chosen deterministically in $0.85$ ms & \textbf{PASS} \\ \hline
\textbf{TC-05} & Signal Safety Phasing & Agent requests immediate phase switch & Enforces min-green (12s), yellow (3s), and all-red (3s) & Red, Yellow, Green lamps sequence cleanly without flutter & \textbf{PASS} \\ \hline
\textbf{TC-06} & Air-Gapped Zero-Cloud Execution & Wi-Fi and Ethernet physically disconnected & Full pipeline executes video detection and signal control offline & Complete video perception and control with zero network calls & \textbf{PASS} \\ \hline
\end{tabularx}
\end{table}

\subsection{Discussion of Test Results}
\begin{itemize}
    \item \textbf{TC-01 \& TC-02 (Perception \& Queue Tracking)}: Verified that the edge vision pipeline reliably isolates auto-rickshaws and two-wheelers, calculating continuous PCU without relying on painted lane boundaries. Stationary queue tracking accurately ignores free-flowing vehicles.
    \item \textbf{TC-03 \& TC-04 (Neural Dynamics \& Policy Inference)}: The Neural World Model achieved a test state MSE of $0.001579$, proving that the surrogate digital twin accurately models physical vehicle arrivals and queue dissipation. The MAPPO Actor generated inference decisions in under $1.0\text{ ms}$, satisfying real-time edge constraints.
    \item \textbf{TC-05 (Safety State Machine)}: Proved that rapid phase fluttering is prevented by the 12-second minimum green constraint. The transition sequence cleanly visits Yellow clearance ($3\text{ s}$), All-Red clearance ($3\text{ s}$), and Green entry.
    \item \textbf{TC-06 (Offline Operation)}: Verified that the system operates with zero external network connectivity, ensuring absolute data privacy and immunity to municipal internet outages.
\end{itemize}

\subsection{Performance Benchmarks}
Table \ref{tab:ch4_performance_metrics} records the latency and throughput of each operational stage measured on the experimental test platform.

\begin{table}[htbp]
\centering
\small
\caption{Measured execution latency and throughput by system component}
\label{tab:ch4_performance_metrics}
\begin{tabularx}{\textwidth}{|p{5.5cm}|p{4.5cm}|X|}
\hline
\textbf{Processing Stage} & \textbf{Measured Benchmark} & \textbf{Operational Condition} \\ \hline
YOLOv8n Video Detection & 28.5 ms per frame (35.1 FPS) & CPU Inference (Intel i5) \\ \hline
YOLOv8n Video Detection (GPU) & 6.2 ms per frame (161 FPS) & NVIDIA RTX 3050 GPU \\ \hline
LAP Centroid Tracking & 1.4 ms per frame & CPU Tracking (50 objects) \\ \hline
Observation Vector Synthesis & 0.35 ms per step & NumPy vectorization (22-dim) \\ \hline
World Model State Step & 0.42 ms per step (2,380 FPS) & PyTorch CPU evaluation \\ \hline
MAPPO Actor Policy Inference & 0.85 ms per decision epoch & PyTorch CPU forward pass \\ \hline
Streamlit Dashboard Frame Refresh & 32.0 ms per frame ($\sim 31$ FPS) & Chrome browser client \\ \hline
\textbf{Total End-to-End Latency} & \textbf{31.5 ms (CPU) / 8.8 ms (GPU)} & \textbf{Exceeds 30 FPS target} \\ \hline
\end{tabularx}
\end{table}

\section{Safety Measures and Fail-Safe Protocols}
\label{sec:ch4_safety_measures}
To guarantee that the autonomous AI system complies with municipal safety regulations, five structural safety mechanisms are enforced at the software layer:
\begin{enumerate}
    \item \textbf{Hard Minimum Green Lock ($T_{\text{min\_green}} = 12\text{ s}$)}: The RL agent is prohibited from switching signals until the active green phase has elapsed for at least 12 seconds, preventing visual driver confusion and dilemma-zone accidents.
    \item \textbf{Mandatory Yellow Change Interval ($T_{\text{yellow}} = 3\text{ s}$)}: Every green termination is followed by a non-negotiable 3-second amber light to clear vehicles traveling at approach speed.
    \item \textbf{Mandatory All-Red Hold Interval ($T_{\text{all\_red}} = 3\text{ s}$)}: An all-red clearance interval holds all competing approaches at red for 3 seconds before granting green to the conflicting phase, ensuring the intersection box is empty.
    \item \textbf{Gridlock Safety Rotation ($T_{\text{max\_green}} = 90\text{ s}$)}: If an agent attempts to hold green on a single approach indefinitely, the controller overrides the RL policy and forces a rotation to prevent cross-street starvation.
    \item \textbf{Fail-Safe Degradation to Fixed-Time Cycle}: If camera failure, video corruption, or lens occlusion occurs, the controller automatically falls back to a pre-programmed 30-second fixed-time signal cycle, ensuring uninterrupted traffic operation.
\end{enumerate}
